% ****** Start of file apssamp.tex ******
%
%   This file is part of the APS files in the REVTeX 4.2 distribution.
%   Version 4.2a of REVTeX, December 2014
%
%   Copyright (c) 2014 The American Physical Society.
%
%   See the REVTeX 4 README file for restrictions and more information.
%
% TeX'ing this file requires that you have AMS-LaTeX 2.0 installed
% as well as the rest of the prerequisites for REVTeX 4.2
%
% See the REVTeX 4 README file
% It also requires running BibTeX. The commands are as follows:
%
%  1)  latex apssamp.tex
%  2)  bibtex apssamp
%  3)  latex apssamp.tex
%  4)  latex apssamp.tex
%
\documentclass[%
 reprint,
%superscriptaddress,
%groupedaddress,
%unsortedaddress,
%runinaddress,
%frontmatterverbose, 
%preprint,
%preprintnumbers,
%nofootinbib,
%nobibnotes,
%bibnotes,
 amsmath,amssymb,
 aps,
pra,
% prb,
%rmp,
%prstab,
%prstper,
%floatfix,
]{revtex4-2}

\usepackage{braket}
\usepackage{hyperref}
\hypersetup{
    colorlinks=true,
    linkcolor=black,
    citecolor=black,
    urlcolor=black,
    }
\usepackage{graphicx}% Include figure files
\usepackage{dcolumn}% Align table columns on decimal point
\usepackage{bm}% bold math
%\usepackage{hyperref}% add hypertext capabilities
%\usepackage[mathlines]{lineno}% Enable numbering of text and display math
%\linenumbers\relax % Commence numbering lines

%\usepackage[showframe,%Uncomment any one of the following lines to test 
%%scale=0.7, marginratio={1:1, 2:3}, ignoreall,% default settings
%%text={7in,10in},centering,
%%margin=1.5in,
%%total={6.5in,8.75in}, top=1.2in, left=0.9in, includefoot,
%%height=10in,a5paper,hmargin={3cm,0.8in},
%]{geometry}

\begin{document}

\title{Fermi--Amaldi-Based Local Mixing Function}% Force line breaks with \\

\author{Ivan P. Bosko and Viktor N. Staroverov}
 % \email{vstarove@uwo.ca}
\affiliation{%
 Department of Chemistry, The University of Western Ontario, London, Ontario N6A 5B7, Canada
}%

\date{\today}

\begin{abstract}

\end{abstract}
\maketitle

\section{Main Text}

Here we propose a new local mixing function (LMF) for local hybrid functionals. This LMF is Fermi--Amaldi (FA) based and represents the ratio between the FA and the exact exchange energy densities. The first is given by 
\begin{equation}
    e_{\text{x}}^{\text{FA}}(\mathbf{r})
    =
    -
    \frac{1}{2N}
    \int\frac{\rho(\mathbf{r}')}{|\mathbf{r}-\mathbf{r}'|}
    d\mathbf{r}',
\end{equation} 
and the latter is expressed via molecular orbitals $\{\varphi_i(\mathbf{r})\}$ and the ground-state density matrix $\mathbf{D}$ as 
\begin{gather}
    e_{\text{x}}^{\text{exact}}(\mathbf{r})\nonumber\\[10pt]
    =
    -\frac{1}{2}
    \sum_{pqrs}
    D_{pq}
    D_{rs}
    \int
    \frac{\varphi^*_p(\mathbf{r})\varphi_s(\mathbf{r})
    \varphi^*_r(\mathbf{r}')\varphi_q(\mathbf{r}')}
    {|\mathbf{r}-\mathbf{r}'|}
    d\mathbf{r}'.
\end{gather}
 
We can show that the proposed LMF satisfies the zero to one boundary range:
\begin{gather*}
    e_{\text{x}}^{\text{exact}}(\mathbf{r})
    \leq
    0,\\[10pt]
    e_{\text{x}}^{\text{FA}}(\mathbf{r})
    \leq
    0,\\[10pt]
    e_{\text{x}}^{\text{exact}}(\mathbf{r})
    \leq
    e_{\text{x}}^{\text{FA}}(\mathbf{r}),
\end{gather*}
thus
\begin{equation}
    0
    \leq
    \frac{e_{\text{x}}^{\text{FA}}}
    {e_{\text{x}}^{\text{exact}}}
    \leq
    1.
\end{equation}

\bibliography{../bibliography}
\end{document}

